Large language models are increasingly used for information, advice, writing, and decision support across culturally different settings. A response can be fluent and factually plausible while still being culturally inappropriate: it may generalize a regional practice, confuse a social norm with a legal rule, overlook religious or historical sensitivities, or represent a group through a stereotype. Evaluating such cases is difficult because culture is contextual and internally diverse, and because not every cultural question has one externally verifiable answer.

This thesis presents \textbf{Vericult}, a standalone, evidence-grounded verifier for the post-hoc evaluation of culturally situated LLM responses. Vericult evaluates one prompt--response pair at a time. It first extracts only context supported by the prompt and selects relevant dimensions from a ten-part cultural rubric. It then identifies up to three material response targets. Externally checkable targets are converted into neutral baseline and variation questions and passed to a candidate-blind evidence stage. The system retrieves and classifies sources, creates an evidence memo, performs at most one follow-up search, and freezes the memo before the original response target is reintroduced. The final stages compare the target with the frozen evidence, score the applicable cultural dimensions, and preserve abstention when the basis is insufficient. All intermediate decisions are stored in structured traces.

The final evaluation contains 360 fixed prompt--response pairs: 120 custom PLT cases, 120 prompts selected from six external multicultural datasets, and 120 externally grounded cultural red-team cases. All responses are generated with the same frozen GPT-OSS 120B generator. Vericult uses a separate Qwen3.6 35B-A3B backbone. Evidence is first acquired in LIVE mode, audited and hashed, and then reused in REPLAY mode. REPLAY is the primary Vericult execution reported in the thesis.

\placeholder{FINAL RESULTS: insert the verified REPLAY completion rate, outcome distribution, evidence coverage, main trace-based finding, and only those baseline comparisons supported by the final data.}

The main contribution is therefore not a claim that web evidence can define a single correct culture. It is a reproducible verification architecture that makes cultural dimensions, external evidence, uncertainty, and failure points more inspectable than a single opaque preference score. The thesis evaluates both the usefulness and the limitations of this approach on the frozen 360-pair dataset.
